% table_relmse.tex
% Table III — Per-tensor reconstruction error (selected tensors, k=16)
%
% Shows where metaheuristics improve over K-means.
% Values: RelMSE of the published weights, results/rev2_logs/regen/
% relmse_published_weights.csv (scripts/relmse_from_bins.py); full per-tensor
% output in results/rev2_logs/regen/final/compress_*.log.
% Selected: tok_embeddings (largest), 3 attention, 3 FFN (representative layers).
%
% RelMSE = MSE / (mean(w^2) + eps)  — relative reconstruction error
%
% Usage: % table_relmse.tex
% Table III — Per-tensor reconstruction error (selected tensors, k=16)
%
% Shows where metaheuristics improve over K-means.
% Values: RelMSE of the published weights, results/rev2_logs/regen/
% relmse_published_weights.csv (scripts/relmse_from_bins.py); full per-tensor
% output in results/rev2_logs/regen/final/compress_*.log.
% Selected: tok_embeddings (largest), 3 attention, 3 FFN (representative layers).
%
% RelMSE = MSE / (mean(w^2) + eps)  — relative reconstruction error
%
% Usage: % table_relmse.tex
% Table III — Per-tensor reconstruction error (selected tensors, k=16)
%
% Shows where metaheuristics improve over K-means.
% Values: RelMSE of the published weights, results/rev2_logs/regen/
% relmse_published_weights.csv (scripts/relmse_from_bins.py); full per-tensor
% output in results/rev2_logs/regen/final/compress_*.log.
% Selected: tok_embeddings (largest), 3 attention, 3 FFN (representative layers).
%
% RelMSE = MSE / (mean(w^2) + eps)  — relative reconstruction error
%
% Usage: % table_relmse.tex
% Table III — Per-tensor reconstruction error (selected tensors, k=16)
%
% Shows where metaheuristics improve over K-means.
% Values: RelMSE of the published weights, results/rev2_logs/regen/
% relmse_published_weights.csv (scripts/relmse_from_bins.py); full per-tensor
% output in results/rev2_logs/regen/final/compress_*.log.
% Selected: tok_embeddings (largest), 3 attention, 3 FFN (representative layers).
%
% RelMSE = MSE / (mean(w^2) + eps)  — relative reconstruction error
%
% Usage: \input{tables/table_relmse}
% Requires: booktabs

\begin{table}[t]
\centering
\footnotesize
\caption{Per-tensor relative reconstruction error (RelMSE) for selected tensors,
         $k{=}16$ centroids. Bold = best per row. The embedding matrix
         (\texttt{tok\_emb}) is the largest tensor (32\,K$\times$288\,=\,9.2\,M weights)
         and the primary contributor to PPL improvement under PSO.}
\label{tab:relmse}
\setlength{\tabcolsep}{4pt}
\renewcommand{\arraystretch}{1.15}
\begin{tabular}{lrrrr}
\toprule
\textbf{Tensor} & \textbf{Shape} & \textbf{K-means} & \textbf{GA} & \textbf{PSO} \\
\midrule
\multicolumn{5}{l}{\textit{Token embedding (largest tensor)}} \\
\quad\texttt{tok\_emb}         & $32000\times288$ & 0.01190 & 0.01191 & \textbf{0.01013} \\
\midrule
\multicolumn{5}{l}{\textit{Attention weights (layer 0)}} \\
\quad\texttt{wq.0}              & $288\times288$   & 0.01553 & 0.01553 & \textbf{0.01564}$^*$ \\
\quad\texttt{wk.0}              & $288\times288$   & 0.01473 & 0.01473 & 0.01474 \\
\quad\texttt{wo.0}              & $288\times288$   & 0.01927 & \textbf{0.01927} & 0.01994 \\
\midrule
\multicolumn{5}{l}{\textit{Feed-forward weights (layer 0)}} \\
\quad\texttt{w1.0}              & $768\times288$   & \textbf{0.00997} & \textbf{0.00997} & 0.00999 \\
\quad\texttt{w2.0}              & $288\times768$   & \textbf{0.01442} & 0.01443 & 0.01447 \\
\quad\texttt{w3.0}              & $768\times288$   & \textbf{0.00968} & \textbf{0.00968} & 0.00969 \\
\midrule
\multicolumn{5}{l}{\textit{Deep layers (layer 5)}} \\
\quad\texttt{wq.5}              & $288\times288$   & \textbf{0.01368} & \textbf{0.01368} & 0.01369 \\
\quad\texttt{w1.5}              & $768\times288$   & \textbf{0.01131} & \textbf{0.01131} & 0.01132 \\
\quad\texttt{w2.5}              & $288\times768$   & \textbf{0.01332} & \textbf{0.01332} & 0.01345 \\
\midrule
\textit{Mean (all 43 tensors)} & — & 0.01264 & 0.01264 & 0.01265 \\
\bottomrule
\end{tabular}
%
\smallskip
\begin{minipage}{0.95\columnwidth}
  \footnotesize
  $^*$~Marginal PSO regressions on small tensors ($\leq$82\,K weights) reflect
  stochastic sampling noise in fitness evaluation (50\,K-sample subsample),
  not genuine degradation. K-means finds near-optimal solutions for small
  tensors, confirmed by the identical mean RelMSE across methods.
  PSO's advantage concentrates on \texttt{tok\_emb}: $-$14.9\,\% RelMSE,
  and PSO lowers held-out TinyStories perplexity from 8.47 (K-means) to 7.79.
\end{minipage}
\end{table}

% Requires: booktabs

\begin{table}[t]
\centering
\footnotesize
\caption{Per-tensor relative reconstruction error (RelMSE) for selected tensors,
         $k{=}16$ centroids. Bold = best per row. The embedding matrix
         (\texttt{tok\_emb}) is the largest tensor (32\,K$\times$288\,=\,9.2\,M weights)
         and the primary contributor to PPL improvement under PSO.}
\label{tab:relmse}
\setlength{\tabcolsep}{4pt}
\renewcommand{\arraystretch}{1.15}
\begin{tabular}{lrrrr}
\toprule
\textbf{Tensor} & \textbf{Shape} & \textbf{K-means} & \textbf{GA} & \textbf{PSO} \\
\midrule
\multicolumn{5}{l}{\textit{Token embedding (largest tensor)}} \\
\quad\texttt{tok\_emb}         & $32000\times288$ & 0.01190 & 0.01191 & \textbf{0.01013} \\
\midrule
\multicolumn{5}{l}{\textit{Attention weights (layer 0)}} \\
\quad\texttt{wq.0}              & $288\times288$   & 0.01553 & 0.01553 & \textbf{0.01564}$^*$ \\
\quad\texttt{wk.0}              & $288\times288$   & 0.01473 & 0.01473 & 0.01474 \\
\quad\texttt{wo.0}              & $288\times288$   & 0.01927 & \textbf{0.01927} & 0.01994 \\
\midrule
\multicolumn{5}{l}{\textit{Feed-forward weights (layer 0)}} \\
\quad\texttt{w1.0}              & $768\times288$   & \textbf{0.00997} & \textbf{0.00997} & 0.00999 \\
\quad\texttt{w2.0}              & $288\times768$   & \textbf{0.01442} & 0.01443 & 0.01447 \\
\quad\texttt{w3.0}              & $768\times288$   & \textbf{0.00968} & \textbf{0.00968} & 0.00969 \\
\midrule
\multicolumn{5}{l}{\textit{Deep layers (layer 5)}} \\
\quad\texttt{wq.5}              & $288\times288$   & \textbf{0.01368} & \textbf{0.01368} & 0.01369 \\
\quad\texttt{w1.5}              & $768\times288$   & \textbf{0.01131} & \textbf{0.01131} & 0.01132 \\
\quad\texttt{w2.5}              & $288\times768$   & \textbf{0.01332} & \textbf{0.01332} & 0.01345 \\
\midrule
\textit{Mean (all 43 tensors)} & — & 0.01264 & 0.01264 & 0.01265 \\
\bottomrule
\end{tabular}
%
\smallskip
\begin{minipage}{0.95\columnwidth}
  \footnotesize
  $^*$~Marginal PSO regressions on small tensors ($\leq$82\,K weights) reflect
  stochastic sampling noise in fitness evaluation (50\,K-sample subsample),
  not genuine degradation. K-means finds near-optimal solutions for small
  tensors, confirmed by the identical mean RelMSE across methods.
  PSO's advantage concentrates on \texttt{tok\_emb}: $-$14.9\,\% RelMSE,
  and PSO lowers held-out TinyStories perplexity from 8.47 (K-means) to 7.79.
\end{minipage}
\end{table}

% Requires: booktabs

\begin{table}[t]
\centering
\footnotesize
\caption{Per-tensor relative reconstruction error (RelMSE) for selected tensors,
         $k{=}16$ centroids. Bold = best per row. The embedding matrix
         (\texttt{tok\_emb}) is the largest tensor (32\,K$\times$288\,=\,9.2\,M weights)
         and the primary contributor to PPL improvement under PSO.}
\label{tab:relmse}
\setlength{\tabcolsep}{4pt}
\renewcommand{\arraystretch}{1.15}
\begin{tabular}{lrrrr}
\toprule
\textbf{Tensor} & \textbf{Shape} & \textbf{K-means} & \textbf{GA} & \textbf{PSO} \\
\midrule
\multicolumn{5}{l}{\textit{Token embedding (largest tensor)}} \\
\quad\texttt{tok\_emb}         & $32000\times288$ & 0.01190 & 0.01191 & \textbf{0.01013} \\
\midrule
\multicolumn{5}{l}{\textit{Attention weights (layer 0)}} \\
\quad\texttt{wq.0}              & $288\times288$   & 0.01553 & 0.01553 & \textbf{0.01564}$^*$ \\
\quad\texttt{wk.0}              & $288\times288$   & 0.01473 & 0.01473 & 0.01474 \\
\quad\texttt{wo.0}              & $288\times288$   & 0.01927 & \textbf{0.01927} & 0.01994 \\
\midrule
\multicolumn{5}{l}{\textit{Feed-forward weights (layer 0)}} \\
\quad\texttt{w1.0}              & $768\times288$   & \textbf{0.00997} & \textbf{0.00997} & 0.00999 \\
\quad\texttt{w2.0}              & $288\times768$   & \textbf{0.01442} & 0.01443 & 0.01447 \\
\quad\texttt{w3.0}              & $768\times288$   & \textbf{0.00968} & \textbf{0.00968} & 0.00969 \\
\midrule
\multicolumn{5}{l}{\textit{Deep layers (layer 5)}} \\
\quad\texttt{wq.5}              & $288\times288$   & \textbf{0.01368} & \textbf{0.01368} & 0.01369 \\
\quad\texttt{w1.5}              & $768\times288$   & \textbf{0.01131} & \textbf{0.01131} & 0.01132 \\
\quad\texttt{w2.5}              & $288\times768$   & \textbf{0.01332} & \textbf{0.01332} & 0.01345 \\
\midrule
\textit{Mean (all 43 tensors)} & — & 0.01264 & 0.01264 & 0.01265 \\
\bottomrule
\end{tabular}
%
\smallskip
\begin{minipage}{0.95\columnwidth}
  \footnotesize
  $^*$~Marginal PSO regressions on small tensors ($\leq$82\,K weights) reflect
  stochastic sampling noise in fitness evaluation (50\,K-sample subsample),
  not genuine degradation. K-means finds near-optimal solutions for small
  tensors, confirmed by the identical mean RelMSE across methods.
  PSO's advantage concentrates on \texttt{tok\_emb}: $-$14.9\,\% RelMSE,
  and PSO lowers held-out TinyStories perplexity from 8.47 (K-means) to 7.79.
\end{minipage}
\end{table}

% Requires: booktabs

\begin{table}[t]
\centering
\footnotesize
\caption{Per-tensor relative reconstruction error (RelMSE) for selected tensors,
         $k{=}16$ centroids. Bold = best per row. The embedding matrix
         (\texttt{tok\_emb}) is the largest tensor (32\,K$\times$288\,=\,9.2\,M weights)
         and the primary contributor to PPL improvement under PSO.}
\label{tab:relmse}
\setlength{\tabcolsep}{4pt}
\renewcommand{\arraystretch}{1.15}
\begin{tabular}{lrrrr}
\toprule
\textbf{Tensor} & \textbf{Shape} & \textbf{K-means} & \textbf{GA} & \textbf{PSO} \\
\midrule
\multicolumn{5}{l}{\textit{Token embedding (largest tensor)}} \\
\quad\texttt{tok\_emb}         & $32000\times288$ & 0.01190 & 0.01191 & \textbf{0.01013} \\
\midrule
\multicolumn{5}{l}{\textit{Attention weights (layer 0)}} \\
\quad\texttt{wq.0}              & $288\times288$   & 0.01553 & 0.01553 & \textbf{0.01564}$^*$ \\
\quad\texttt{wk.0}              & $288\times288$   & 0.01473 & 0.01473 & 0.01474 \\
\quad\texttt{wo.0}              & $288\times288$   & 0.01927 & \textbf{0.01927} & 0.01994 \\
\midrule
\multicolumn{5}{l}{\textit{Feed-forward weights (layer 0)}} \\
\quad\texttt{w1.0}              & $768\times288$   & \textbf{0.00997} & \textbf{0.00997} & 0.00999 \\
\quad\texttt{w2.0}              & $288\times768$   & \textbf{0.01442} & 0.01443 & 0.01447 \\
\quad\texttt{w3.0}              & $768\times288$   & \textbf{0.00968} & \textbf{0.00968} & 0.00969 \\
\midrule
\multicolumn{5}{l}{\textit{Deep layers (layer 5)}} \\
\quad\texttt{wq.5}              & $288\times288$   & \textbf{0.01368} & \textbf{0.01368} & 0.01369 \\
\quad\texttt{w1.5}              & $768\times288$   & \textbf{0.01131} & \textbf{0.01131} & 0.01132 \\
\quad\texttt{w2.5}              & $288\times768$   & \textbf{0.01332} & \textbf{0.01332} & 0.01345 \\
\midrule
\textit{Mean (all 43 tensors)} & — & 0.01264 & 0.01264 & 0.01265 \\
\bottomrule
\end{tabular}
%
\smallskip
\begin{minipage}{0.95\columnwidth}
  \footnotesize
  $^*$~Marginal PSO regressions on small tensors ($\leq$82\,K weights) reflect
  stochastic sampling noise in fitness evaluation (50\,K-sample subsample),
  not genuine degradation. K-means finds near-optimal solutions for small
  tensors, confirmed by the identical mean RelMSE across methods.
  PSO's advantage concentrates on \texttt{tok\_emb}: $-$14.9\,\% RelMSE,
  and PSO lowers held-out TinyStories perplexity from 8.47 (K-means) to 7.79.
\end{minipage}
\end{table}
